\documentclass[a4paper]{article}
\usepackage{ctex}
\usepackage[affil-it]{authblk}
\usepackage[backend=bibtex,style=numeric]{biblatex}
\usepackage{amsmath,amsthm,amssymb,amsfonts}
\usepackage{graphicx}
\usepackage{bm}
\usepackage{float}
\usepackage{geometry}
\geometry{margin=1.5cm, vmargin={0pt,1cm}}
\setlength{\topmargin}{-1cm}
\setlength{\paperheight}{29.7cm}
\setlength{\textheight}{25.3cm}

\addbibresource{citation.bib}

\begin{document}
% =================================================
\title{微分方程数值解 2026 春夏编程作业三}

\author{曾申昊 3220100701
  \thanks{邮箱: \texttt{3097714673@qq.com}
                            \texttt{或} \texttt{3220100701@zju.edu.cn}}}
\affil{浙江大学}

\date{更新时间: \today}

\maketitle

% ============================================

\section{理论分析}

\par  一个 s 步的线性多步法的通用形式为
\begin{equation}
  \sum_{j=0}^s \alpha_j \mathbf{U}^{n+j} 
  = k \sum_{j=0}^s \beta_j \mathbf{f}(\mathbf{U}^{n+j}, t_{n+j})
\end{equation}
\par 其中 $\alpha_j$ 和 $\beta_j$ 是常数，且 $\alpha_s \neq 0$。
讲义上已经给出了 ABM、AMM 和 BDF 的 $\alpha_j$ 和 $\beta_j$ 的值。
\par 对于这三类方法，在具体实现时，都需要使用单步方法来计算前 $s$ 步的数值解，以提供初始条件，
具体程序实现时，选择了经典的 RK4 方法来计算前 $s$ 步的数值解。
\par 另外，对于ABM ，由于它是显式方法，因此在每一步的计算中，
数值解 $\mathbf{U}^{n+s}$ 可以直接通过前 $s$ 步计算得到。
而对于 AMM 和 BDF ，由于它们是隐式方法，因此在每一步的计算中，
需要通过求解一个非线性方程来得到数值解 $\mathbf{U}^{n+s}$。
在具体实现时，选择了不动点迭代法来求解这个非线性方程，初始值设置为 RK4 方法计算得到的数值解。

\section{数值结果}

% ===============================================

\par 下面是 Adams-Bashforth 方法在 Case 2 中的数值结果，
表格中列出了不同步数 $s$ 和阶数 $p$ 的数值解的误差和 CPU 时间。
可以看到，随着步长 $k$ 的减小，各阶方法的误差按照相应的阶数进行收敛，
同时 CPU 时间也随着步长的减小而成倍增加。
这里用 Case 2 来展示数值结果，是因为相对来说 Case 1 的数值解需要更小的步长才能达到较好的精度。
对于 Adams-Moulton 方法和 BDF 方法上述结论同样适用，除了
5阶的 Adams-Moulton 方法在 Case 2 中只有4阶精度，
可能是由于初始条件是用 RK4 方法计算得到的。
\par 另外还测试了三种方法在 Case 1 和 Case 2
中至少需要多大的步长才能使图形解能够较好地拟合。 

\begin{table}
    \centering
    \caption{Adams-Bashforth 方法 Case 2 的数值结果}
    \begin{tabular}{|c|c|c|c|c|c|c|}
        \hline
        步数$s$ & 阶数$p$ & 步长$k$ & $L^1$ 误差  & $L^2$ 误差 & $L^\infty$ 误差 & CPU时间 \\
        \hline
        1 & 1 & 0.00004 & 0.7549788151 & 0.5963579678 & 0.5852949421 & 4.976192 \\
        1 & 1 & 0.00002 & 0.3983907089 & 0.2759142661 & 0.2627883653 & 9.854486 \\
        1 & 1 & 0.00001 & 0.1362794474 & 0.09231627652 & 0.08794388448 & 19.370515 \\
        1 & 1 & 0.000005 & 0.06174725987 & 0.03852857422 & 0.03502963777 & 38.398576 \\
        \hline 
        2 & 3 & 0.0008 & 0.001221709517 & 0.0007935709459 & 0.0007193317996 & 0.35911 \\
        2 & 3 & 0.0004 & 0.0003582208294 & 0.000230712439 & 0.0002093800426 & 0.613865 \\
        2 & 3 & 0.0002 & 0.0001163189891 & 6.412550308e-05 & 4.841096853e-05 & 1.229305 \\
        2 & 3 & 0.0001 & 2.393770298e-05 & 1.558380915e-05 & 1.413556898e-05 & 2.377439 \\
        \hline 
        3 & 4 & 0.0032 & 0.03211062092 & 0.01886681694 & 0.01605825298 & 0.104099 \\
        3 & 4 & 0.0016 & 0.004368573107 & 0.002474602319 & 0.002045126585 & 0.207764 \\
        3 & 4 & 0.0008 & 0.0006429019013 & 0.0003506679515 & 0.0002380177996 & 0.3728 \\
        3 & 4 & 0.0004 & 6.477002824e-05 & 3.844610882e-05 & 3.319558205e-05 & 0.759138 \\
        \hline 
        4 & 5 & 0.016 & 0.2500349211 & 0.1389341579 & 0.1068189898 & 0.023527 \\ 
        4 & 5 & 0.008 & 0.01103506582 & 0.006014372013 & 0.004486401465 & 0.046369 \\
        4 & 5 & 0.004 & 0.0002670413642 & 0.0001517950463 & 0.0001257085727 & 0.090162 \\
        4 & 5 & 0.002 & 5.406020787e-06 & 2.922110307e-06 & 2.056200311e-06 & 0.182866 \\
        \hline 
    \end{tabular}
\end{table}

\begin{table}
    \centering
    \caption{Adams-Moulton 方法 Case 2 的数值结果}
    \begin{tabular}{|c|c|c|c|c|c|c|}
        \hline
        步数$s$ & 阶数$p$ & 步长$k$ & $L^1$ 误差  & $L^2$ 误差 & $L^\infty$ 误差 & CPU时间 \\
        \hline 
        1 & 2 & 0.0032 & 0.004982849005 & 0.003259907949 & 0.002960476103 & 1.329572 \\
        1 & 2 & 0.0016 & 0.001597706004 & 0.0009278825496 & 0.0007908525933 & 2.650371 \\
        1 & 2 & 0.0008 & 0.0004902468968 & 0.0002806197883 & 0.0002117485899 & 5.176032 \\
        1 & 2 & 0.0004 & 7.820594179e-05 & 5.108376824e-05 & 4.632609437e-05 & 10.20849 \\
        \hline 
        2 & 3 & 0.0128 & 0.2826057837 & 0.1894406035 & 0.1797245155 & 0.452687 \\
        2 & 3 & 0.0064 & 0.02877872784 & 0.01778716783 & 0.01570743125 & 0.90283 \\
        2 & 3 & 0.0032 & 0.003684144816 & 0.002187862064 & 0.001889655675 & 1.77264 \\
        2 & 3 & 0.0016 & 0.000702144342 & 0.0004000427828 & 0.0002885501614 & 3.55934 \\
        \hline
        3 & 4 & 0.032 & 0.5360752022 & 0.400014759 & 0.3893349184 & 0.243016 \\ 
        3 & 4 & 0.016 & 0.01227439088 & 0.007159137595 & 0.006096933911 & 0.459361 \\
        3 & 4 & 0.008 & 0.001959685142 & 0.001522029605 & 0.001461337247 & 0.920643 \\
        3 & 4 & 0.004 & 6.776283353e-06 & 3.666098967e-06 & 2.608377872e-06 & 1.797678 \\
        \hline
        4 & 5 & 0.032 & 0.226829394 & 0.1287005739 & 0.1030204775 & 0.268655 \\ 
        4 & 5 & 0.016 & 0.01143972754 & 0.006772321694 & 0.005816396897 & 0.572801 \\
        4 & 5 & 0.008 & 0.001820278075 & 0.001573686647 & 0.001565406192 & 1.100188 \\
        4 & 5 & 0.004 & 1.248278796e-05 & 7.4203724e-06 & 6.41822728e-06 & 2.260583 \\
        \hline 
    \end{tabular}
\end{table}

\begin{table}
    \centering
    \caption{Backward Difference Formula 方法 Case 2 的数值结果}
    \begin{tabular}{|c|c|c|c|c|c|c|}
        \hline
        步数$s$ & 阶数$p$ & 步长$k$ & $L^1$ 误差  & $L^2$ 误差 & $L^\infty$ 误差 & CPU时间 \\
        \hline
        1 & 1 & 0.00008 & 0.3133657914 & 0.1596014923 & 0.09510297345 & 34.287378 \\
        1 & 1 & 0.00004 & 0.2286501334 & 0.1201746863 & 0.07621336922 & 68.377356 \\
        1 & 1 & 0.00002 & 0.150807538 & 0.08196152141 & 0.05866780353 & 135.855435 \\
        \hline 
        2 & 2 & 0.0008 & 0.001164997344 & 0.0006826089285 & 0.0005900224298 & 3.603495 \\
        2 & 2 & 0.0004 & 0.0002800597611 & 0.0001797425144 & 0.0001631624534 & 7.167146 \\
        2 & 2 & 0.0002 & 0.0001187391428 & 6.815043583e-05 & 5.025845892e-05 & 14.196881 \\
        2 & 2 & 0.0001 & 1.905031643e-05 & 1.239199431e-05 & 1.124117225e-05 & 28.24953 \\
        \hline 
        3 & 3 & 0.0032 & 0.02243938411 & 0.01340760966 & 0.01166883279 & 0.938906 \\
        3 & 3 & 0.0016 & 0.002506804772 & 0.001650676212 & 0.001468184414 & 1.899695 \\
        3 & 3 & 0.0008 & 0.000439688062 & 0.0002498919195 & 0.0002036555917 & 3.852212 \\
        3 & 3 & 0.0004 & 4.318335163e-05 & 2.56336539e-05 & 2.21339179e-05 & 7.538902 \\
        \hline
        4 & 4 & 0.016 & 0.2364630071 & 0.1512205735 & 0.1402633484 & 0.200332 \\ 
        4 & 4 & 0.008 & 0.006149711166 & 0.003753521841 & 0.003315531434 & 0.410071 \\
        4 & 4 & 0.004 & 0.0001683182159 & 9.603371133e-05 & 7.992624886e-05 & 0.821637 \\
        4 & 4 & 0.002 & 3.574528051e-06 & 1.940888489e-06 & 1.422934978e-06 & 1.575101 \\
        \hline 
    \end{tabular}
\end{table}

\begin{table}
    \centering
    \caption{三种方法数值解拟合的最小步长（ABM, AMM, BDF）}
    \begin{tabular}{|c|c|c|}
        \hline
        $p$ & Case 1 步长 & Case 2 步长 \\
        \hline
        1 & 0.00001 & 0.0001 \\
        2 & 0.0001 & 0.01 \\
        3 & 0.0006 & 0.01 \\
        4 & 0.0007 & 0.02 \\
        \hline 
    \end{tabular}
    \begin{tabular}{|c|c|c|}
        \hline
        $p$ & Case 1 步长 & Case 2 步长 \\
        \hline
        2 & 0.0003 & 0.02 \\
        3 & 0.001 & 0.02 \\
        4 & 0.001 & 0.03 \\
        5 & 0.001 & 0.07 \\
        \hline 
    \end{tabular}
    \begin{tabular}{|c|c|c|}
        \hline
        $p$ & Case 1 步长 & Case 2 步长 \\
        \hline
        1 & 0.00002 & 0.0001 \\
        2 & 0.0002 & 0.01 \\
        3 & 0.0005 & 0.01 \\
        4 & 0.001 & 0.02 \\
        \hline 
    \end{tabular}
\end{table}

\printbibliography

\end{document}